% Live Verify — arXiv-track paper. Build with `make` (tectonic).
% Draft: sections marked [PLANNED] await their studies; do not submit before removing them
% or completing the work they describe.
\documentclass[11pt]{article}

\usepackage[margin=1.1in]{geometry}
\usepackage{booktabs}
\usepackage{longtable}
\usepackage{amssymb}
\usepackage{microtype}
\usepackage[hidelinks]{hyperref}

% GENERATED by paper/scripts/gen-conformance-table.js — do not hand-edit.
\newcommand{\VectorTotal}{17}
\newcommand{\TextVectorCount}{8}
\newcommand{\ImageVectorCount}{9}
\newcommand{\TextVectorPasses}{8}
\newcommand{\CharNormVectorCount}{1}
\newcommand{\OcrRuleVectorCount}{0}


\title{Live Verify: Verifying Printed and Displayed Text\\
Against Issuer Domains by Canonicalized Hashing\thanks{Draft. Reference
implementations, the normative specification, and the conformance corpus are at
\url{https://github.com/live-verify/live-verify}.}}
\author{Paul Hammant}
\date{Draft of \today}

\begin{document}
\maketitle

\begin{abstract}
Documents that circulate as paper, PDF, or on-screen text routinely outlive any
channel that could vouch for them: a determination letter is photographed, a
license is framed on a wall, a bank statement is forwarded. Live Verify is a
deliberately low-fidelity bridge from such artifacts back to the party that
issued them. A document carries a human-readable \texttt{verify:} line naming
the issuer's domain; a verifier's client canonicalizes the visible text,
hashes it with SHA-256, and performs an HTTPS GET for that hash at the issuer's
domain. A 200 response with an affirming status means the issuer, \emph{now},
stands behind exactly that text. The design trades cryptographic evidence for
radical deployability: there are no signatures, no key ceremonies, and nothing
to install on the document side beyond one printable line. We specify the
canonicalization pipeline that makes camera-captured and clipboard-captured
text hash identically across implementations, publish a conformance corpus of
hash-pinned vectors that three implementations pass, and give an honest
account of the trust model's limits: verification is a live confirmation, not
transferable proof; it is deniable by deletion; and public hash endpoints act
as confirmation oracles for guessable documents, which constrains the class of
documents the scheme should carry.
\end{abstract}

\section{Introduction}

Institutions already know how to publish facts on their own domains, and the
web's public-key infrastructure already binds those domains to organizational
identity well enough for banking. What is missing is a bridge from a
\emph{document in hand} --- paper, screenshot, forwarded PDF --- to that
issuer-controlled surface. Existing bridges either carry opaque payloads (QR
codes, NFC chips) that a human cannot audit and a printer cannot always
reproduce, or demand heavyweight adoption (digital signatures, verifiable
credentials) from issuers whose documents are still fundamentally text.

Live Verify's premise is that the document's own visible text can be the
payload. The scheme adds one line:

\begin{center}
\texttt{verify:example.com/path}
\end{center}

A verifying client captures the text above that line by camera or clipboard,
canonicalizes it (\S\ref{sec:canon}), hashes it, and asks the named domain
whether it stands behind that hash. The answer is rendered against the
\emph{authority domain} --- the domain printed on the document --- never
against whichever host happened to serve bytes.

\paragraph{Contributions.}
\begin{enumerate}
  \item A protocol for text-first document verification whose entire
        issuer-side deployment is static file hosting
        (\S\ref{sec:design}).
  \item A canonicalization pipeline reconciling OCR capture and digital
        clipboard capture into one hash space, with the order of operations
        specified normatively (\S\ref{sec:canon}).
  \item A published, append-only conformance corpus in which each vector's
        filename \emph{is} its pinned SHA-256, and cross-implementation
        results over it (\S\ref{sec:eval}).
  \item An explicit negative result offered as a design position: what a
        live-lookup scheme \emph{cannot} provide (non-repudiation,
        transferable proof, resistance to deletion), and why naming those
        limits is a precondition for deploying it safely
        (\S\ref{sec:limits}).
\end{enumerate}

\section{Related Work}\label{sec:related}

\paragraph{Domain-anchored transparency.} Certificate Transparency
\cite{rfc6962} demonstrated that append-only, publicly auditable logs can
discipline an existing trust ecosystem without redesigning it. Live Verify
shares the instinct of anchoring trust in operational domain control rather
than new key material, but deliberately omits the log: the issuer's live
endpoint is the (revocable, deniable) source of truth, with third-party
witnessing left as an optional layer.

\paragraph{Hash-prefix privacy.} Have I Been Pwned's $k$-anonymity range
queries \cite{ali2018hibp} showed that a hash-lookup service can avoid
learning which exact record was queried. Live Verify's baseline lookup is a
full-hash GET and therefore leaks the queried hash to the issuer; adopting
prefix range queries is designed but unimplemented
(\S\ref{sec:limits}).

\paragraph{Verifiable credentials.} The W3C VC model \cite{w3cvc} carries
cryptographic proof inside the credential, achieving non-repudiation and
offline verification at the cost of issuer-side signing infrastructure and
holder-side wallets. Live Verify occupies the opposite corner of the design
space: nothing travels with the document but text.

\paragraph{Provenance grading.} The Isn\=ad--Rij\=al framework
\cite{raja2026isnad} grades chains of transmission in multi-agent knowledge
systems, importing classical hadith methodology. Its concerns are
complementary: Live Verify authenticates a leaf artifact against its issuer,
while isn\=ad-style grading evaluates the chain through which derived
knowledge traveled. The overlap and its limits are examined in the
repository's comparative documents.

\section{Protocol Design}\label{sec:design}

\subsection{The verify line}
The verify line is the last content line of a claim, in ABNF \cite{rfc5234}:
\begin{verbatim}
verify-line = scheme ":" verify-target
scheme      = "verify" / "vfy"
\end{verbatim}
Recognition is tolerant (case-insensitive, whitespace around the colon,
bottom-up scan) because the line must survive OCR. A verify line immediately
preceded by \texttt{>} is \emph{quoted} --- a claim reproduced inside another
claim --- and is inert; OCR artifact cleanup deliberately does not strip
\texttt{>} so that quotation survives the camera path.

\subsection{Canonicalization}\label{sec:canon}
Camera capture and clipboard capture of the same document must reach the same
bytes. A conforming implementation applies, in order: (1) OCR artifact
cleanup, camera path only --- stripping border characters and scan debris;
(2) Unicode canonical composition (NFC) \cite{uax15}, since precomposed and
decomposed encodings of the same glyph otherwise hash differently; (3)
issuer-declared single-character folds; (4) issuer-declared OCR regex rules,
treated as untrusted input and guarded; (5) a fixed punctuation fold table
(curly quotes, en/em dashes, no-break space, ellipsis); (6) line processing
--- trim, collapse internal whitespace, drop blank lines, rejoin with LF, no
trailing newline. The hash is SHA-256 \cite{fips180} over the UTF-8 encoding
of the result, rendered as 64 lowercase hex characters.

\subsection{Lookup and affirmation}
The client GETs \texttt{https://host[/path]/\{hash\}} --- HTTPS
unconditionally, with a parsed-hostname (never substring) exception for
\texttt{localhost} development. A claim is verified iff the response is HTTP
200 and its status affirms, either as JSON \texttt{status:"verified"} or via
an issuer-declared response-type whose class is \texttt{affirming}. Responses
never echo document content: the verifier already holds the document, and an
endpoint that repeats content is a data leak.

\subsection{Endorsement}
An issuer's metadata may name an endorser. The endorsement commits to a hash
of the issuer's \emph{entire} metadata file, so any change to the issuer's
self-description invalidates it; clients walk at most three links and render
unendorsed domains as self-verified with an explicit caution.

\section{Trust Model}\label{sec:trust}

Verification here means: \emph{the issuer's domain, at this moment, affirms
this exact text}. That is a live confirmation with the strength --- and only
the strength --- of WebPKI domain authentication plus the issuer's
operational integrity. The repository's security document treats each threat
in turn (query privacy, transport and DNS posture, issuer-supplied
normalization, spoofing); \S\ref{sec:limits} summarizes what survives
honest scrutiny.

\section{Evaluation}\label{sec:eval}

\subsection{Cross-implementation hash determinism}
The conformance corpus currently contains \VectorTotal{} vectors:
\TextVectorCount{} text vectors (each filename a pinned SHA-256) and
\ImageVectorCount{} image vectors exercising the full camera pipeline on
Android (ML~Kit) and iOS (Vision). The reference JavaScript implementation
passes \TextVectorPasses/\TextVectorCount{} text vectors; the table in
Appendix~\ref{app:vectors} is \emph{generated by executing the implementation
against the corpus at paper build time}, so the numbers cannot drift from the
code. Vectors are discriminating by design --- the NFC vector's body is
decomposed Unicode, and an implementation that skips canonical composition
fails it.

\subsection{OCR field failure study [PLANNED]}
Which characters do production OCR stacks actually confuse on real documents
(e-ink, laser print, screens), and does the issuer-fold mechanism cover the
observed failure classes? This study is planned; anecdotal evidence from
development (e.g.\ e-ink capture) motivates it but is not evidence.

\subsection{Confirmation-oracle entropy analysis [PLANNED]}
For which document classes is the plaintext guessable enough that a public
hash endpoint becomes a confirmation oracle? A quantified entropy analysis
across the use-case corpus is planned; the security document currently
handles this qualitatively (high-entropy fields, optional salt).

\section{Limitations}\label{sec:limits}

Live Verify's verification is not evidence. There is no signature, hence no
non-repudiation: an issuer can delete a hash and deny ever having affirmed
it. Confirmation is not transferable: a verifier cannot prove to a third
party that verification succeeded. The full-hash lookup leaks query patterns
to issuers; $k$-anonymity range queries \cite{ali2018hibp} and OPRF-based
lookups are designed but not implemented. An HTTP 404 is both ``retry after
network trouble'' and ``the issuer denies this document''; treating a
security signal as retryable trains verifiers past it, and client design must
separate the two. These properties bound the scheme's honest use: it
complements, and does not replace, signature-bearing schemes where evidence
must outlive the issuer's cooperation.

\section{Conclusion}

Live Verify makes one narrow bet: for documents that are already text, the
cheapest verification bridge that institutions will actually deploy is a
printable line, a canonicalization discipline, and static file hosting on a
domain they already defend. The specification, three reference
implementations, and a hash-pinned conformance corpus are published; the two
planned studies above are the empirical work that remains.

\bibliographystyle{plain}
\bibliography{refs}

\appendix
\section{Conformance vectors (generated)}\label{app:vectors}
% GENERATED by paper/scripts/gen-conformance-table.js — do not hand-edit.
\begin{longtable}{p{8.2cm} l c}
\toprule
\textbf{Vector (description)} & \textbf{Pinned hash (prefix)} & \textbf{JS} \\
\midrule
\endhead
\bottomrule
\endfoot
Bachelor of Thaumatology training page certificate & \texttt{1cddfbb2adfa} & \checkmark \\
Curly double quotes normalized to straight quotes & \texttt{30c6d7a77f52} & \checkmark \\
NFC canonicalization - body uses decomposed Unicode (base letters + combining accents), curly quotes, ellipsis and em-dash; implementations must apply NFC then the standard folds before hashing, or this vector fails & \texttt{4970a05e322a} & \checkmark \\
En-dash and em-dash normalized to hyphen & \texttt{50093137b7cc} & \checkmark \\
Master of Applied Anthropics training page certificate & \texttt{6725b845dcdf} & \checkmark \\
Simple text - baseline test & \texttt{a591a6d40bf4} & \checkmark \\
Document-specific character normalization (diacritics) & \texttt{ea9837801631} & \checkmark \\
Whitespace normalization - leading, trailing, multiple spaces, blank lines & \texttt{eb4108396033} & \checkmark \\
\end{longtable}


\end{document}
